\documentclass[11pt]{article}
\usepackage{amsmath,amssymb}
\usepackage[margin=25mm]{geometry}
\title{SNSPD Model Equations}
\author{}
\date{}
\begin{document}
\maketitle

\section*{Equations shown in the presentation}
% Slide 1: the two explicit resistive RF branches.
\[
50\,\Omega\parallel50\,\Omega=25\,\Omega .
\]
This is the parallel resistance of the two RF shunts. The complete
frequency-dependent detector load also includes the series return and
reactive elements.

% Slides 2--4: non-inverting gain and nominal gain blocks.
\[
K=1+\frac{R_{\mathrm{f}}}{R_{\mathrm{g}}},
\qquad G=12.281 .
\]
\[
\text{Section 1: }R_{\mathrm{f}}=15.2\,\mathrm{k}\Omega,
\qquad
\text{Section 2: }R_{\mathrm{f}}=123.5\,\mathrm{k}\Omega,
\qquad R_{\mathrm{g}}=100\,\mathrm{k}\Omega .
\]

% Slide 5: conceptual electrical equation.
\[
L_k\frac{\mathrm{d}I}{\mathrm{d}t}
=V_{\mathrm{d}}-V_{\mathrm{s}}-I R_{\mathrm{wire}} .
\]
\[
R_{\mathrm{wire}}
=\min\!\left[R_{\mathrm{n}},
\max\!\left(R_{\mathrm{seed}}+R_{\mathrm{d}},0\right)\right] .
\]
Here $R_{\mathrm{seed}}$ is the heated-seed resistance,
$R_{\mathrm{d}}$ is the propagated-domain resistance, and
$R_{\mathrm{n}}$ is the fully normal wire resistance.
The conceptual circuit equation omits the numerical series resistance
$R_{\mathrm{s,num}}=10\,\mu\Omega$. The implemented electrical equation
includes the additional term $-I R_{\mathrm{s,num}}$.

% Slide 6: annotated internal structure; V = I R_wire is shown.
% Auxiliary capacitors implement numerical state integrators.

% Slide 7: seed energy balance. These are total seed heat capacities.
\[
C_{\mathrm{e}}\frac{\mathrm{d}T_{\mathrm{e}}}{\mathrm{d}t}
=P_{\mathrm{opt}}+I^2R_{\mathrm{seed}}
-P_{\mathrm{e-ph}}-P_{\mathrm{diff}} .
\]
\[
C_{\mathrm{ph}}\frac{\mathrm{d}T_{\mathrm{ph}}}{\mathrm{d}t}
=P_{\mathrm{e-ph}}-P_{\mathrm{sub}} .
\]
The quantities $C_{\mathrm{e}}$ and $C_{\mathrm{ph}}$ are total heat
capacities of the seed, not heat capacities per unit volume.
The electron--phonon power enters the two balances with opposite signs.

% Slide 8: Joule-heating label in the feedback diagram.
\[
I^2R_{\mathrm{seed}} .
\]
The propagated-domain diagnostic $I^2R_{\mathrm{d}}$ is not added to
the seed heat balance in this model.

% Slide 9: integrated optical power over the deposited pulse.
\[
E_{\mathrm{dep}}=\int P_{\mathrm{opt}}(t)\,\mathrm{d}t .
\]
The integration is over the deposited optical pulse.

\section*{Heat-transfer laws in the model notes}
\[
P_{\mathrm{e-ph}}
=\Sigma V_{\mathrm{seed}}
\left(T_{\mathrm{e}}^5-T_{\mathrm{ph}}^5\right),
\qquad
P_{\mathrm{sub}}
=\beta A_{\mathrm{seed}}
\left(T_{\mathrm{ph}}^4-T_{\mathrm{b}}^4\right) .
\]
The seed volume is $V_{\mathrm{seed}}$, its substrate-contact area is
$A_{\mathrm{seed}}$, and the bath temperature is $T_{\mathrm{b}}$.

\section*{Notation and source}
Descriptive subscripts are upright. The differential symbol is written
as $\mathrm{d}$, while variables are italic. Units are upright and
separated from numbers by a thin space. The built-in operators
$\min$ and $\max$ are upright.

The equations were checked against the following files in the SNSPD
simulation package:
\begin{itemize}
\item \texttt{models/snspd\_nbn\_v5c.lib}
\item \texttt{circuits/bias\_tee\_and\_load.inc}
\item \texttt{circuits/amplifier\_chain.inc}
\end{itemize}
This is a notation and presentation check; no model parameter or
simulation equation was changed.
\end{document}
